\section{What the macro describes: the spec types}
\label{sec:spec}

\subsection{The types (\texttt{core/teeny-core/src/model/tile\_spec.rs})}

\begin{lstlisting}
KernelTileSpec   { inputs: &[TensorTileSpec], outputs: &[TensorTileSpec], loop_spec: Option<TileLoopSpec> }
TensorTileSpec   { param, rank, axes: &[TileAxisBinding], reduction_axis: Option<usize>, untiled_dims: &[&str] }
TileAxisBinding  { dims: &[usize], block_const, extent_param, window: Option<TileWindow>, divide_by: Option<usize> }
TileWindow       { output_extent_param, stride_const, pad_const: Option<&str>, kernel_size_const }
TileLoopSpec     { carries: &[TileCarryBinding], trip_count_factors: &[&str] }
TileCarryBinding { name, shape_consts: &[&str] }            // entries: const names or "1"
\end{lstlisting}

There are still no enums and no variants. Every ``case'' is a combination of field
values. Every string is a \emph{name}, never a value: the spec has no per-node
parameter lookup, so a consumer that needs numbers must be given a
\code{ConstLookup} by its caller. \code{KernelTileSpec::validate()} checks the
invariants the types cannot hold:
\begin{itemize}
\item dims are below the rank and appear at most once;
\item no name is both bound and untiled;
\item there are no more described dims than the rank;
\item every carry has a shape;
\item there is at least one output, unless the spec is loop-only.
\end{itemize}

\subsection{How a declaration becomes a binding}

The same per-axis rule applies on both routes:

{\small
\begin{xltabular}{\linewidth}{>{\raggedright\arraybackslash}p{4.8cm}X}
\toprule
\textbf{Declared axis} & \textbf{Emitted} \\
\midrule
\endhead
\code{block = B} (not \code{walk}) & A \code{TileAxisBinding { dims, block_const: "B", extent_param, window }} \\
\code{block = B, walk} & \emph{No binding}; its name goes to \code{untiled_dims}. One program walks the whole axis, so it cannot be split across programs. A binding would tell the scheduler it can be. \\
no \code{block}, with \code{window(..)} & A binding with the literal \code{block_const: "1"}, so the window has something to hang off. \\
no \code{block}, no window & No binding; the name goes to \code{untiled_dims}. \\
\code{reduce} & Sets \code{reduction_axis} to the axis's real dim. \\
\code{span = N} & \code{dims} is the next \code{N} real dims; the rank grows by $N-1$. \\
\bottomrule
\end{xltabular}
}

Which parameters appear, and with what rank:
\begin{itemize}
\item \textbf{Route 1, several declared axes:} \code{tile_spec()} at a \emph{fixed}
  rank. Each parameter reports its own axes, so a \code{(C,)} bias against an
  \code{[N, C]} output is rank~1 (\issue{1tl.4}).
\item \textbf{Route 1, one flat axis:} \code{tile_spec(rank)}, with one binding over
  \code{dims = [0..rank)}. The same flat kernel applies at any rank, so the graph node
  supplies it. The lowering calls \code{SiluForward::<f32>::tile_spec(node.shape.len())}.
\item \textbf{Route 2:} \code{tile_spec()} at fixed rank. \code{InOut} parameters
  appear in both lists.
\item \textbf{Loop only:} flash attention declares its loop and no tensors
  (\code{BH} is a grid extent, not a parameter), so it gets a spec with empty
  \code{inputs}/\code{outputs}. \code{validate()} accepts it; \code{resolve_inputs}
  rejects it.
\item \textbf{An input with no \code{#[tile]} is omitted from the spec entirely.}
  conv2d's weight is the important case: its second extent is \code{C_IN / G},
  which \code{extent = <ident>} cannot name. It is read through
  \code{#[tile_loop_scalar]} and does not appear in the spec.
\end{itemize}

\subsection{GridSpec and FusionCore}

\code{grid_spec()} lists the output's non-walked axes as \code{GridAxisBinding { name,
extent_factors, dim, block_const }}, in dim order. If a \code{#[tile_grid(order =
..)]} is present it uses that order and subset instead. For example:
\begin{itemize}
\item \code{transpose_2d_forward} decodes \code{M} outer while its output is
  \code{[N, M]}.
\item \code{batch_norm_normalize_forward}'s grid is \code{[C]} alone.
\item \code{conv2d_bn_silu_tiled} decodes \code{(B, OH, C_OUT, OW)}.
\end{itemize}
\code{swizzled} emits no grid spec at all. The purpose is the ``rider'' check of
\issue{teenygrad-1nr.17}: a fused op can reuse the anchor's already-decoded indices
only if its grid axes are a name-matching subset of the anchor's.

\code{fusion_core()} returns \code{Some(FusionCore { input_ident, output_ident,
body_source, extra_params })} (the body minus its trailing store) only when all of
these hold:
\begin{itemize}
\item exactly one \code{In<Tile>} and one \code{Out<Tile>};
\item each declares one axis;
\item the hand-written body ends in \code{T::store(..)};
\item there is no \code{generate} loop and no loop-indexed operand.
\end{itemize}
\code{silu_forward} qualifies. \code{batch_norm_2d_nchw_forward_inference} (five
inputs) and \code{conv2d_forward} (generated loop) do not. Whether a generated-loop
kernel even \emph{has} a fusion core is recorded as undecided in
\issue{teenygrad-1cli}. \S\ref{sec:welder-infer} proposes an answer.
